\section{Perturbation Methods}
Moradi and Samwald (2021) introduce a broad range of input perturbations at two granular levels – \textbf{character-level} and \textbf{word-level} – to simulate realistic noise in text ([Evaluating the Robustness of Neural Language Models to Input Perturbations](https://aclanthology.org/2021.emnlp-main.117.pdf#:~:text=Insertion%20Who%20was%20the%20first,How%20many%20heartts%20does%20an)) ([Evaluating the Robustness of Neural Language Models to Input Perturbations](https://aclanthology.org/2021.emnlp-main.117.pdf#:~:text=Negation%20What%20planet%20is%20known,the%20common%20most%20color%20eye)). These perturbations include minor typographical errors and linguistic tweaks that a human might accidentally or intentionally produce, without drastically altering the overall meaning of the input. Below we summarize the perturbation types and their specific forms:

- \textbf{Character-Level Perturbations:} Small edits to individual characters in a word. Seven methods were implemented:  
  - \textit{Insertion:} Inserting an extra random character into a word (e.g., “first” → \textbf{“firsdt”}) ([Evaluating the Robustness of Neural Language Models to Input Perturbations](https://aclanthology.org/2021.emnlp-main.117.pdf#:~:text=Insertion%20Who%20was%20the%20first,How%20many%20heartts%20does%20an)).  
  - \textit{Deletion:} Removing a character from a word (e.g., “discovered” → \textbf{“discovred”}) ([Evaluating the Robustness of Neural Language Models to Input Perturbations](https://aclanthology.org/2021.emnlp-main.117.pdf#:~:text=Insertion%20Who%20was%20the%20first,How%20many%20heartts%20does%20an)).  
  - \textit{Replacement:} Substituting a character with a different one, often simulating a typo (e.g., “Minister” → \textbf{“Monister”}) ([Evaluating the Robustness of Neural Language Models to Input Perturbations](https://aclanthology.org/2021.emnlp-main.117.pdf#:~:text=Replacement%20Who%20is%20the%20Prime,How%20many%20heartts%20does%20an)).  
  - \textit{Swapping:} Exchanging the positions of two adjacent characters (e.g., “language” → \textbf{“lnaguage”}) ([Evaluating the Robustness of Neural Language Models to Input Perturbations](https://aclanthology.org/2021.emnlp-main.117.pdf#:~:text=Replacement%20Who%20is%20the%20Prime,How%20many%20heartts%20does%20an)).  
  - \textit{Character Repetition:} Repeating a character twice in a row (e.g., “hearts” → \textbf{“heartts”}) ([Evaluating the Robustness of Neural Language Models to Input Perturbations](https://aclanthology.org/2021.emnlp-main.117.pdf#:~:text=Swapping%20What%20is%20the%20primary,hearts%20does%20an%20OCTOPUS%20have)).  
  - \textit{Common Misspelling (CMW):} Replacing a word with a commonly misspelled variant from a predefined list (e.g., “fluorescent” → \textbf{“florescent”}) ([Evaluating the Robustness of Neural Language Models to Input Perturbations](https://aclanthology.org/2021.emnlp-main.117.pdf#:~:text=CMW%20What%20kind%20of%20gas,hearts%20does%20an%20OCTOPUS%20have)).  
  - \textit{Letter Case Change (LCC):} Altering letter casing, such as lowercasing or uppercasing characters (e.g., “octopus” → \textbf{“OCTOPUS”}) ([Evaluating the Robustness of Neural Language Models to Input Perturbations](https://aclanthology.org/2021.emnlp-main.117.pdf#:~:text=CMW%20What%20kind%20of%20gas,rock%20is%20a%20form%20of)).  
  These character-level perturbations preserve most letters of the original word, creating realistic typos or format changes that a human reader can usually still understand. They expose models to variation in spelling and casing (for example, simulating a user typing error or inconsistent capitalization).

- \textbf{Word-Level Perturbations:} Edits that involve whole words, targeting lexical or syntactic properties of the input. Seven methods were implemented:  
  - \textit{Word Deletion:} Dropping a random word from the input sentence (e.g., “How much \textbf{was} a ticket for the Titanic?” → “How much a ticket for the Titanic?”) ([Evaluating the Robustness of Neural Language Models to Input Perturbations](https://aclanthology.org/2021.emnlp-main.117.pdf#:~:text=Deletion%20How%20much%20was%20a,known%20as%20the%20%E2%80%9Cred%E2%80%9D%20planet)). This simulates missing words or telegraphic speech.  
  - \textit{Word Repetition:} Duplicating a word in place, i.e. inserting an extra copy of an existing word (e.g., “What is another \textbf{name} for vitamin B1?” → “What is another \textbf{name name} for vitamin B1?”) ([Evaluating the Robustness of Neural Language Models to Input Perturbations](https://aclanthology.org/2021.emnlp-main.117.pdf#:~:text=Deletion%20How%20much%20was%20a,known%20as%20the%20%E2%80%9Cred%E2%80%9D%20planet)). This mimics stutters or emphasis.  
  - \textit{Replacement with Synonym:} Substituting a word with one of its synonyms from WordNet ([Evaluating the Robustness of Neural Language Models to Input Perturbations](https://aclanthology.org/2021.emnlp-main.117.pdf#:~:text=Repetition%20What%20is%20another%20name,Why%20in%20tennis%20were%20zero)). For example, “precious stone” might become “\textbf{valued rock\textbf{” in a question ([Evaluating the Robustness of Neural Language Models to Input Perturbations](https://aclanthology.org/2021.emnlp-main.117.pdf#:~:text=Repetition%20What%20is%20another%20name,Why%20in%20tennis%20were%20zero)). This tests semantic invariance under lexical choice changes.  
  - \textit{Negation Injection:} Flipping the polarity of a verb phrase by adding or removing negation ([Evaluating the Robustness of Neural Language Models to Input Perturbations](https://aclanthology.org/2021.emnlp-main.117.pdf#:~:text=Negation%20What%20planet%20is%20known,the%20common%20most%20color%20eye)). For instance, “What planet \textbf{is known} as the 'red' planet?” → “What planet \textbf{is not known} as the 'red' planet?” ([Evaluating the Robustness of Neural Language Models to Input Perturbations](https://aclanthology.org/2021.emnlp-main.117.pdf#:~:text=RWS%20What%20precious%20stone%20is,known%20as%20the%20%E2%80%9Cred%E2%80%9D%20planet)). This perturbation changes the sentence’s meaning by introducing negation, challenging the model’s understanding of logical negation and its effect on the task.  
  - \textit{Singular/Plural Verb Form (SPV):} Altering verb number agreement, e.g. using a plural verb form in place of singular or vice versa ([Evaluating the Robustness of Neural Language Models to Input Perturbations](https://aclanthology.org/2021.emnlp-main.117.pdf#:~:text=Negation%20What%20planet%20is%20known,Why%20in%20tennis%20were%20zero)). For example, “What \textbf{does} a barometer measure?” → “What \textbf{do} a barometer measure?” ([Evaluating the Robustness of Neural Language Models to Input Perturbations](https://aclanthology.org/2021.emnlp-main.117.pdf#:~:text=Negation%20What%20planet%20is%20known,Why%20in%20tennis%20were%20zero)). This creates a grammatically incorrect sentence (subject-verb disagreement) without changing the core intent of the question.  
  - \textit{Verb Tense Change:} Changing the tense of verbs in the sentence (present ↔ past, etc.) ([Evaluating the Robustness of Neural Language Models to Input Perturbations](https://aclanthology.org/2021.emnlp-main.117.pdf#:~:text=Negation%20What%20planet%20is%20known,were%20zero%20points%20called%20love)). For example, “Why in tennis \textbf{are} zero points called love?” → “Why in tennis \textbf{were} zero points called love?” ([Evaluating the Robustness of Neural Language Models to Input Perturbations](https://aclanthology.org/2021.emnlp-main.117.pdf#:~:text=Negation%20What%20planet%20is%20known,were%20zero%20points%20called%20love)). This can introduce subtle meaning shifts or temporal framing differences while largely preserving the topic of the sentence.  
  - \textit{Word Order Shuffling:} Reordering a selection of words in the sentence out of their original sequence ([Evaluating the Robustness of Neural Language Models to Input Perturbations](https://aclanthology.org/2021.emnlp-main.117.pdf#:~:text=Negation%20What%20planet%20is%20known,the%20common%20most%20color%20eye)). For instance, “What is the \textbf{most common eye color\textbf{?” → “What is the \textbf{common most color eye\textbf{?” ([Evaluating the Robustness of Neural Language Models to Input Perturbations](https://aclanthology.org/2021.emnlp-main.117.pdf#:~:text=Negation%20What%20planet%20is%20known,the%20common%20most%20color%20eye)). This jumbles the syntax (often yielding an awkward or mildly ungrammatical sentence) but retains the same set of words/content.  
  
Each perturbation is designed to be \textbf{automatically generated}. The authors relied on simple rules or lexical resources (like WordNet for synonyms, or a list of common misspellings) and tools (POS taggers for negation insertion, etc.) to implement these transformations ([Evaluating the Robustness of Neural Language Models to Input Perturbations](https://aclanthology.org/2021.emnlp-main.117.pdf#:~:text=4.2%20Word,from%20the%20WordNet%20lexical%20database)) ([Evaluating the Robustness of Neural Language Models to Input Perturbations](https://aclanthology.org/2021.emnlp-main.117.pdf#:~:text=Negation,the%20firsdt%20governor%20of%20Alaska)). Notably, these perturbations are \textit{non-adversarial} in nature – they were not crafted to fool the model in a malicious way, but rather to mimic natural noise (typos, grammatical variations, etc.) that could occur in real-world user input. By applying such perturbations systematically, the study evaluates how robust state-of-the-art language models are when faced with inputs that deviate slightly from the training distribution ([Evaluating the Robustness of Neural Language Models to Input Perturbations - ACL Anthology](https://aclanthology.org/2021.emnlp-main.117/#:~:text=handling%20different%20types%20of%20input,understanding%20of%20NLP%20systems%E2%80%99%20robustness)) ([Evaluating the Robustness of Neural Language Models to Input Perturbations](https://aclanthology.org/2021.emnlp-main.117.pdf#:~:text=only%20applied%20to%20the%20test,specific%20types%20of%20perturbation%20more)).

\section{Safety Mitigation Techniques}
To ensure the robustness evaluation was reliable and fair, the authors employed several mitigation and validation techniques as part of their methodology:

- \textbf{Human Evaluation of Perturbations (User Study):} An extensive user study was conducted to verify that the perturbed inputs remained \textit{understandable} and, in most cases, preserved the original task-relevant meaning ([Evaluating the Robustness of Neural Language Models to Input Perturbations](https://aclanthology.org/2021.emnlp-main.117.pdf#:~:text=In%20the%20first%20part%20of,every%20perturbed%20sample%2C%20and%20were)) ([Evaluating the Robustness of Neural Language Models to Input Perturbations](https://aclanthology.org/2021.emnlp-main.117.pdf#:~:text=In%20the%20second%20part%20of,if)). Twenty participants were asked to judge perturbed examples side-by-side with the original text. This evaluation was split into two parts: (1) perturbations assumed to preserve meaning (all character-level and certain word-level perturbations like repetition, singular/plural change, and verb tense change), and (2) perturbations likely to alter meaning (word-level deletion, synonym replacement, negation, and word order changes) ([Evaluating the Robustness of Neural Language Models to Input Perturbations](https://aclanthology.org/2021.emnlp-main.117.pdf#:~:text=In%20the%20first%20part%20of,every%20perturbed%20sample%2C%20and%20were)) ([Evaluating the Robustness of Neural Language Models to Input Perturbations](https://aclanthology.org/2021.emnlp-main.117.pdf#:~:text=methods%20that%20may%20change%20the,with%20the%20test%20set%20label)). Results showed that \textbf{≈94%} of samples with meaning-preserving perturbations were indeed judged understandable and semantically equivalent to the original ([Evaluating the Robustness of Neural Language Models to Input Perturbations](https://aclanthology.org/2021.emnlp-main.117.pdf#:~:text=According%20to%20the%20user%20evaluations%2C,to%20ensure%20understandability%20and%20consistency)). In contrast, for perturbations that can change meaning, only \textbf{≈39%} of samples retained their original label/meaning, with the remainder either needing label changes or becoming nonsensical ([Evaluating the Robustness of Neural Language Models to Input Perturbations](https://aclanthology.org/2021.emnlp-main.117.pdf#:~:text=According%20to%20the%20user%20evaluations%2C,on%20the%20respective%20original%20inputs)). This human filtering process ensured that the final test sets used for robustness evaluation contained mostly valid perturbations – i.e. the input is perturbed but still a legitimate instance of the task so that any performance drop reflects model brittleness rather than a genuine label change ([Evaluating the Robustness of Neural Language Models to Input Perturbations](https://aclanthology.org/2021.emnlp-main.117.pdf#:~:text=94,of%20the%20study%2C%20each%20participant)) ([Evaluating the Robustness of Neural Language Models to Input Perturbations](https://aclanthology.org/2021.emnlp-main.117.pdf#:~:text=According%20to%20the%20user%20evaluations%2C,three%20NLP%20tasks%20for%20which)). Essentially, the user study served as a quality control, helping the authors curate perturbed test inputs that are \textit{fair} for the models to be evaluated on (especially for tricky cases like negation or deletion). Only perturbations passing this human check were kept for analysis, thereby mitigating the risk of mis-evaluating a model on nonsense or semantically flipped inputs.

- \textbf{Model Variety and Comparison:} The study evaluated four different neural language models – BERT, RoBERTa, XLNet (all Transformer-based models), and ELMo (a biLSTM with character convolutions) – under the same perturbation scenarios ([Evaluating the Robustness of Neural Language Models to Input Perturbations](https://aclanthology.org/2021.emnlp-main.117.pdf#:~:text=RoBERTa%20still%20performs%20better%20than,be%20due%20to%20its%20pure)) ([Evaluating the Robustness of Neural Language Models to Input Perturbations](https://aclanthology.org/2021.emnlp-main.117.pdf#:~:text=results%20also%20suggest%20those%20models,based%20model%2C%20i.e.%20ELMo%2C%20is)). By comparing multiple architectures, the authors identified whether some model types are inherently more robust than others. This acts as a diagnostic tool: if certain models consistently handle perturbations better, it suggests potential avenues for safer model designs or training regimes. For example, including ELMo (which uses character-level representations) allowed them to observe the impact of having char-level information on robustness, as compared to BERT/RoBERTa’s subword tokenization ([Evaluating the Robustness of Neural Language Models to Input Perturbations](https://aclanthology.org/2021.emnlp-main.117.pdf#:~:text=can%20handle%20specific%20types%20of,to%20handle%20perturbations%20to%20word)). Such cross-model evaluation is a mitigation strategy in itself – it highlights vulnerabilities and strengths of each model, informing how future systems might be improved for robustness. The study found notable differences (discussed below), indicating that architecture and pretraining can significantly affect a model’s tolerance to input noise ([Evaluating the Robustness of Neural Language Models to Input Perturbations](https://aclanthology.org/2021.emnlp-main.117.pdf#:~:text=RoBERTa%20still%20performs%20better%20than,be%20due%20to%20its%20pure)) ([Evaluating the Robustness of Neural Language Models to Input Perturbations](https://aclanthology.org/2021.emnlp-main.117.pdf#:~:text=XLNet%20is%20shown%20to%20handle,robust%20when%20words%20are%20replaced)).

- \textbf{Robustness Benchmarking and Guidelines:} Based on their findings, the authors advocate for incorporating \textbf{perturbation-based evaluation} into standard model assessment. They highlight that conventional benchmarks (with clean, unperturbed text) often paint an overly optimistic picture of model performance, and propose that \textbf{testing on perturbed inputs should become routine} to better reflect real-world reliability ([Evaluating the Robustness of Neural Language Models to Input Perturbations - ACL Anthology](https://aclanthology.org/2021.emnlp-main.117/#:~:text=handling%20different%20types%20of%20input,understanding%20of%20NLP%20systems%E2%80%99%20robustness)) ([Evaluating the Robustness of Neural Language Models to Input Perturbations](https://aclanthology.org/2021.emnlp-main.117.pdf#:~:text=only%20applied%20to%20the%20test,specific%20types%20of%20perturbation%20more)). In essence, this is a recommendation for a \textit{mitigation at the evaluation level} – ensuring model safety and robustness by routinely checking how models handle slight input variations. By doing so, developers and researchers can catch brittle behavior early and address it (for instance, via data augmentation or improved training methods). The paper stops short of implementing a specific defense (it focuses on evaluation), but this recommendation is a key take-away for the community: robust and safe NLP systems need \textit{perturbation-aware} evaluation and possibly training. Additionally, the perturbation framework developed (with code released) can be used alongside tools like CheckList or adversarial testing frameworks to automatically generate noisy test samples, reducing the burden on users to craft such tests ([Evaluating the Robustness of Neural Language Models to Input Perturbations](https://aclanthology.org/2021.emnlp-main.117.pdf#:~:text=models%20on%20various%20NLP%20tasks,we%20used%20five%20datasets%20covering)) ([Evaluating the Robustness of Neural Language Models to Input Perturbations](https://aclanthology.org/2021.emnlp-main.117.pdf#:~:text=only%20applied%20to%20the%20test,specific%20types%20of%20perturbation%20more)). The authors also note that most of their perturbation methods (except those like negation or deletion which need monitoring) can be applied automatically to create meaningful variations, providing a practical toolkit for robustness testing ([Evaluating the Robustness of Neural Language Models to Input Perturbations](https://aclanthology.org/2021.emnlp-main.117.pdf#:~:text=94,of%20the%20study%2C%20each%20participant)) ([Evaluating the Robustness of Neural Language Models to Input Perturbations](https://aclanthology.org/2021.emnlp-main.117.pdf#:~:text=methods%20that%20may%20change%20the,with%20the%20test%20set%20label)).

In summary, the study’s methodology was designed with careful checks (via human evaluation) to ensure \textit{safety} and validity of perturbations, and it compares multiple models and recommends broader evaluation practices as mitigation steps. These efforts help separate true robustness issues from trivial artifacts and set a precedent for evaluating NLP models under noisy conditions as a standard part of their safety assessment.

\section{Evaluation and Analysis}
\textbf{Overall Sensitivity:} The experiments revealed that even \textit{small input perturbations can lead to significant performance degradation} in state-of-the-art language models ([Evaluating the Robustness of Neural Language Models to Input Perturbations](https://aclanthology.org/2021.emnlp-main.117.pdf#:~:text=only%20applied%20to%20the%20test,specific%20types%20of%20perturbation%20more)) ([Evaluating the Robustness of Neural Language Models to Input Perturbations](https://aclanthology.org/2021.emnlp-main.117.pdf#:~:text=performance%20can%20decrease%20even%20when,current%20benchmarks%20are%20not%20reflecting)). On unperturbed (original) test sets, the models achieved high accuracy/F1 scores, but introducing minor noise (typos, extra/repeated words, etc.) often caused a noticeable drop in those metrics. In many cases the model’s prediction changed from a correct answer to an incorrect one due solely to a trivial alteration of the input. This underscores a reliability concern: models that appear highly accurate in clean benchmarks may still be brittle when faced with the kind of minor errors humans routinely overlook. The authors emphasize that current benchmarks do \textbf{not adequately reflect robustness}, and they argue for complementing them with perturbed-input evaluations for a more realistic gauge of model performance ([Evaluating the Robustness of Neural Language Models to Input Perturbations - ACL Anthology](https://aclanthology.org/2021.emnlp-main.117/#:~:text=handling%20different%20types%20of%20input,understanding%20of%20NLP%20systems%E2%80%99%20robustness)) ([Evaluating the Robustness of Neural Language Models to Input Perturbations](https://aclanthology.org/2021.emnlp-main.117.pdf#:~:text=only%20applied%20to%20the%20test,specific%20types%20of%20perturbation%20more)). In short, all tested models were \textit{sensitive to input noise}, confirming that robustness is a common weakness in NLP systems if not explicitly addressed.

\textbf{Performance Degradation Patterns:} The vulnerability is more pronounced for certain perturbation types and tasks. Notably, the models were \textbf{more sensitive to character-level perturbations than to word-level perturbations on average} ([Evaluating the Robustness of Neural Language Models to Input Perturbations](https://aclanthology.org/2021.emnlp-main.117.pdf#:~:text=model%2C%20the%20average%20of%20absolute,the%20other%20hand%2C%20the%20question)) ([Evaluating the Robustness of Neural Language Models to Input Perturbations](https://aclanthology.org/2021.emnlp-main.117.pdf#:~:text=perturbations%20for%20every%20NLP%20task,other%20tasks%20from%20noisy%20inputs)). Character-level typos can confuse the subword tokenizers or character encoders, leading to unseen or mis-segmented tokens. For example, a single-letter swap or deletion could reduce a model’s accuracy by a large margin. The study quantitatively showed larger drops in performance (accuracy, F1, etc.) under character attacks across all tasks, compared to synonym replacements or grammatical tweaks ([Evaluating the Robustness of Neural Language Models to Input Perturbations](https://aclanthology.org/2021.emnlp-main.117.pdf#:~:text=perturbations%20for%20every%20NLP%20task,other%20tasks%20from%20noisy%20inputs)). This is intuitive since a random character error often yields an OOV (out-of-vocabulary) token or forces the model to rely on sparse character-level features, whereas a word-level change might still share semantic or subword overlap with the original text. That said, certain \textbf{word-level perturbations that do alter semantics (like negation or important word deletion)} can be just as disruptive. In particular, adding or removing a negation in a critical place flipped the meaning and frequently led models to wrong outputs in tasks like sentiment analysis ([Evaluating the Robustness of Neural Language Models to Input Perturbations](https://aclanthology.org/2021.emnlp-main.117.pdf#:~:text=XLNet%20are%20more%20robust%20when,the%20results%2C%20we%20can%20also)). The authors point out that when a perturbation \textit{directly impacts the task’s core signal}, models struggle significantly ([Evaluating the Robustness of Neural Language Models to Input Perturbations](https://aclanthology.org/2021.emnlp-main.117.pdf#:~:text=XLNet%20are%20more%20robust%20when,based%20models)). For instance, inserting “not” in a positive sentiment sentence makes it negative – models often failed to appropriately switch their prediction, yielding errors or inconsistencies.

\textbf{Model-wise Robustness:} Among the models tested, there were clear differences in robustness: \textbf{RoBERTa} showed the best overall stability to perturbations, \textbf{ELMo} the worst, with \textbf{BERT} and \textbf{XLNet} in between ([Evaluating the Robustness of Neural Language Models to Input Perturbations](https://aclanthology.org/2021.emnlp-main.117.pdf#:~:text=perturbations%20and%20their%20performance%20decreases,obtains%20the%20lowest%20scores%20in)). RoBERTa’s superior performance is attributed to its larger pre-training corpus and training optimizations, which apparently gave it a slight edge in handling out-of-distribution inputs ([Evaluating the Robustness of Neural Language Models to Input Perturbations](https://aclanthology.org/2021.emnlp-main.117.pdf#:~:text=results%20also%20suggest%20those%20models,based%20model%2C%20i.e.%20ELMo%2C%20is)) ([Evaluating the Robustness of Neural Language Models to Input Perturbations](https://aclanthology.org/2021.emnlp-main.117.pdf#:~:text=XLNet%20is%20shown%20to%20handle,robust%20when%20words%20are%20replaced)). ELMo, on the other hand, consistently underperformed on perturbed inputs (even though it uses character-level representations) – for example, under certain noise conditions ELMo’s accuracy/F1 dropped so much that it scored far below the Transformer models ([Evaluating the Robustness of Neural Language Models to Input Perturbations](https://aclanthology.org/2021.emnlp-main.117.pdf#:~:text=perturbations%20and%20their%20performance%20decreases,obtains%20the%20lowest%20scores%20in)). Interestingly, ELMo did \textit{outperform BERT on some character-level perturbations} ([Evaluating the Robustness of Neural Language Models to Input Perturbations](https://aclanthology.org/2021.emnlp-main.117.pdf#:~:text=can%20handle%20specific%20types%20of,to%20handle%20perturbations%20to%20word)). The authors note that ELMo’s purely character-based word representations can leverage morphological clues, so it handled some typos better than BERT (which relies on fixed word-piece vocabulary) ([Evaluating the Robustness of Neural Language Models to Input Perturbations](https://aclanthology.org/2021.emnlp-main.117.pdf#:~:text=can%20handle%20specific%20types%20of,to%20handle%20perturbations%20to%20word)). This was seen in cases like minor spelling mistakes, where ELMo could still form a reasonable representation from char convnets, whereas BERT might produce an [UNK] or fragmented token ([Evaluating the Robustness of Neural Language Models to Input Perturbations](https://aclanthology.org/2021.emnlp-main.117.pdf#:~:text=can%20handle%20specific%20types%20of,to%20handle%20perturbations%20to%20word)). However, ELMo’s overall lower capacity and LSTM-based architecture made it \textit{more fragile to word-order changes} and other syntax-level perturbations ([Evaluating the Robustness of Neural Language Models to Input Perturbations](https://aclanthology.org/2021.emnlp-main.117.pdf#:~:text=XLNet%20are%20more%20robust%20when,based%20models)). For example, when words were shuffled, ELMo’s performance degraded more than the self-attention models, since Transformers (BERT, RoBERTa, XLNet) can somewhat cope with reordering due to their permutation-invariant attention mechanism, while LSTMs are inherently order-sensitive ([Evaluating the Robustness of Neural Language Models to Input Perturbations](https://aclanthology.org/2021.emnlp-main.117.pdf#:~:text=XLNet%20is%20shown%20to%20handle,robust%20when%20words%20are%20replaced)) ([Evaluating the Robustness of Neural Language Models to Input Perturbations](https://aclanthology.org/2021.emnlp-main.117.pdf#:~:text=stable%20and%20handle%20the%20noise,based%20models)). In fact, \textbf{XLNet} (which was trained with a permutation language modeling objective) handled jumbled word order best among the models ([Evaluating the Robustness of Neural Language Models to Input Perturbations](https://aclanthology.org/2021.emnlp-main.117.pdf#:~:text=model%20to%20use%20morphological%20clues%2C,suggest%20those%20models%20that%20were)), confirming that its pretraining helps resilience to order permutations. Likewise, models with bigger pretraining datasets (RoBERTa, XLNet) coped better with synonyms substitution than BERT ([Evaluating the Robustness of Neural Language Models to Input Perturbations](https://aclanthology.org/2021.emnlp-main.117.pdf#:~:text=results%20also%20suggest%20those%20models,based%20model%2C%20i.e.%20ELMo%2C%20is)), likely because they had seen more vocabulary variations during training. These model-wise insights suggest that \textit{pretraining scale and objective, as well as architecture, affect robustness}: larger, more diverse pretraining improves tolerance to lexical variation, and char-level modeling helps with typos, whereas autoregressive/permutation training helps with word order issues. No model was immune to perturbations, but RoBERTa was consistently strongest (often achieving the highest perturbed accuracy on many tasks) ([Evaluating the Robustness of Neural Language Models to Input Perturbations](https://aclanthology.org/2021.emnlp-main.117.pdf#:~:text=perturbations%20and%20their%20performance%20decreases,obtains%20the%20lowest%20scores%20in)), and \textit{remains the model to beat in this robustness evaluation}. The fact that even RoBERTa saw substantial drops indicates plenty of room for improvement in future robust model training.

\textbf{Task-wise Sensitivity:} The impact of perturbations also varied across NLP tasks. The study covered a diverse set of tasks (see \textbf{Datasets} below), and found that \textbf{some tasks are inherently more vulnerable to input noise}. In particular, \textbf{Sentiment Analysis (SA)} had the largest performance drops under perturbations ([Evaluating the Robustness of Neural Language Models to Input Perturbations](https://aclanthology.org/2021.emnlp-main.117.pdf#:~:text=negation%20perturbation%20has%20more%20impact,the%20language%20models%20for%20different)) ([Evaluating the Robustness of Neural Language Models to Input Perturbations](https://aclanthology.org/2021.emnlp-main.117.pdf#:~:text=perturbations%20for%20every%20NLP%20task,other%20tasks%20from%20noisy%20inputs)). Small wording changes in sentiment classification (especially negations or intensity changes) often led the models to misclassify the sentiment. For example, inserting an irrelevant word or a typo might distract the model from a key sentiment word, and adding “not” can flip polarity which many models failed to handle correctly ([Evaluating the Robustness of Neural Language Models to Input Perturbations](https://aclanthology.org/2021.emnlp-main.117.pdf#:~:text=XLNet%20are%20more%20robust%20when,based%20models)). This suggests sentiment models rely on specific cue words and are not robust when those cues are altered or when the input is slightly noisy – a serious issue if deployed in user-generated text environments (where typos and slang are common). On the other hand, \textbf{Question Answering (QA)} was found to be relatively \textbf{least affected} by perturbations ([Evaluating the Robustness of Neural Language Models to Input Perturbations](https://aclanthology.org/2021.emnlp-main.117.pdf#:~:text=perturbations%20for%20every%20NLP%20task,other%20tasks%20from%20noisy%20inputs)). The WikiQA task (answer selection) showed smaller accuracy decreases, implying that the QA model could often still find the correct answer sentence despite input noise. A possible reason is that the key information in QA might be redundant or less sensitive to slight phrasing changes – e.g., a misspelling might still match sufficiently with correct answer text, or the model might leverage context around a corrupted token. Additionally, QA metrics may reward partial correctness, and the model might ignore unimportant function words (so dropping or repeating a function word doesn’t always derail comprehension in QA). The authors explicitly note that \textbf{models made more errors in the sentiment analysis task than in other tasks when faced with perturbations, whereas the QA task suffered the least} ([Evaluating the Robustness of Neural Language Models to Input Perturbations](https://aclanthology.org/2021.emnlp-main.117.pdf#:~:text=perturbations%20for%20every%20NLP%20task,other%20tasks%20from%20noisy%20inputs)). The other tasks – \textbf{text classification} (TREC question type classification), \textbf{NER}, and \textbf{semantic similarity} – fell somewhere in between. For instance, \textbf{NER (Named Entity Recognition)} was very sensitive to character-level noise: a misspelling or casing change in a named entity often caused the tagger to fail to recognize it, leading to big F1 drops (the model might label a slightly misspelled “Canada” as O (non-entity) instead of LOCATION). This aligns with intuition: NER models heavily rely on exact spelling (and capitalization features) for entity detection, so they struggled particularly with the CMW and LCC perturbations (the paper’s data showed NER F1 plunging by tens of points in those cases). \textbf{Semantic textual similarity (STS)} also suffered when important content words were perturbed – e.g. swapping or replacing words could lower the similarity score even if a human would judge the two sentences as still similar ([Evaluating the Robustness of Neural Language Models to Input Perturbations](https://aclanthology.org/2021.emnlp-main.117.pdf#:~:text=Figure%201%3A%20Six%20examples%20of,on%20the%20respective%20original%20inputs)). In general, tasks that depend on specific keywords or precise token identity (NER, classification) are undermined by char-level errors, whereas tasks that rely on broader sentence meaning (STS, QA) might sometimes withstand a typo if the rest of the sentence provides enough clue. Still, every task experienced some performance erosion, underscoring that \textit{no NLP task is fully robust to input perturbations}. 

\textbf{Quantitative Highlights:} Across all perturbation types and tasks, the authors observed non-trivial drops in performance metrics. For example, on the SST sentiment task, BERT’s accuracy fell by over 20 points under certain perturbations (compared to the original accuracy) ([Evaluating the Robustness of Neural Language Models to Input Perturbations](https://aclanthology.org/2021.emnlp-main.117.pdf#:~:text=perturbations%20for%20every%20NLP%20task,other%20tasks%20from%20noisy%20inputs)). On the CoNLL NER task, RoBERTa’s F1 dropped from ~92 to ~65 when evaluated on inputs with common misspellings – a dramatic degradation showing the model failing to generalize to slight spelling errors. ELMo’s performance on some perturbed sets was nearly \textit{half} of its original (especially for case changes, where it struggled acutely) ([Evaluating the Robustness of Neural Language Models to Input Perturbations](https://aclanthology.org/2021.emnlp-main.117.pdf#:~:text=CMW%20What%20kind%20of%20gas,rock%20is%20a%20form%20of)) ([Evaluating the Robustness of Neural Language Models to Input Perturbations](https://aclanthology.org/2021.emnlp-main.117.pdf#:~:text=perturbations%20and%20their%20performance%20decreases,obtains%20the%20lowest%20scores%20in)). By contrast, the smallest drops observed were in the WikiQA task with word-level perturbations, where a model like XLNet might lose only ~5-10 points in accuracy for moderate noise levels ([Evaluating the Robustness of Neural Language Models to Input Perturbations](https://aclanthology.org/2021.emnlp-main.117.pdf#:~:text=perturbations%20for%20every%20NLP%20task,other%20tasks%20from%20noisy%20inputs)). To quantify robustness, the paper also computed an “absolute decrease” in performance for each model, perturbation, and task (summarized in their Table 5) ([Evaluating the Robustness of Neural Language Models to Input Perturbations](https://aclanthology.org/2021.emnlp-main.117.pdf#:~:text=Table%205%20presents%20the%20absolute,Perturbed%20inputs%20causes%20the)) ([Evaluating the Robustness of Neural Language Models to Input Perturbations](https://aclanthology.org/2021.emnlp-main.117.pdf#:~:text=sentiment%20analysis%20task%20more%20often,other%20tasks%20from%20noisy%20inputs)). Those results confirm that \textbf{character-level perturbations typically cause higher error rate increases than word-level ones} for all models, and that \textbf{SA is the hardest hit task while QA is the most robust} ([Evaluating the Robustness of Neural Language Models to Input Perturbations](https://aclanthology.org/2021.emnlp-main.117.pdf#:~:text=perturbations%20for%20every%20NLP%20task,other%20tasks%20from%20noisy%20inputs)). They also show that increasing the number of perturbations per sample (applying multiple noises to the same input) exacerbates performance drops in roughly linear fashion ([Evaluating the Robustness of Neural Language Models to Input Perturbations](https://aclanthology.org/2021.emnlp-main.117.pdf#:~:text=Character,8)) – e.g., with four perturbations in a sentence, accuracy dips much more severely than with just one. This indicates compounding effects when noise accumulates. 

In summary, the evaluation vividly demonstrated the fragility of high-performing NLP models: \textit{small input perturbations often led to disproportionate performance declines}. Certain models (RoBERTa, XLNet) and tasks (QA) showed a bit more resilience, whereas others (ELMo, sentiment analysis) were especially brittle. The qualitative takeaway is that current LMs lack robust generalization to slight input variance – a safety concern for real-world deployment. The authors conclude that robustness must be treated as a first-class criterion going forward, suggesting that model training and testing should explicitly address these perturbation scenarios ([Evaluating the Robustness of Neural Language Models to Input Perturbations - ACL Anthology](https://aclanthology.org/2021.emnlp-main.117/#:~:text=handling%20different%20types%20of%20input,understanding%20of%20NLP%20systems%E2%80%99%20robustness)) ([Evaluating the Robustness of Neural Language Models to Input Perturbations](https://aclanthology.org/2021.emnlp-main.117.pdf#:~:text=only%20applied%20to%20the%20test,specific%20types%20of%20perturbation%20more)).

\section{Datasets}
To thoroughly evaluate robustness, the paper uses \textbf{five different NLP datasets} spanning a range of tasks ([Evaluating the Robustness of Neural Language Models to Input Perturbations](https://aclanthology.org/2021.emnlp-main.117.pdf#:~:text=2%20NLP%20tasks%20In%20our,is%20given%20in%20the%20following)) ([Evaluating the Robustness of Neural Language Models to Input Perturbations](https://aclanthology.org/2021.emnlp-main.117.pdf#:~:text=3%20Language%20models%20In%20our,We%20used)). Each dataset corresponds to a distinct task type, allowing the study to assess whether robustness issues are consistent across tasks or more severe in some contexts. Below is an overview of each dataset and its task:

- \textbf{TREC (Question Classification):} A dataset of open-domain factoid questions derived from the TREC QA benchmarks. It contains over \textbf{6,000 questions} (e.g., \textit{“What is the capital of Canada?”}) categorized into \textbf{50 fine-grained question types} ([Evaluating the Robustness of Neural Language Models to Input Perturbations](https://aclanthology.org/2021.emnlp-main.117.pdf#:~:text=TREC%20,specify%20the%20type%20of%20questions)) (such as \textit{Location}, \textit{Person}, \textit{Numeric Information}, etc.). The task is Text Classification (TC): given a question, the model must predict its category. Performance is typically measured by classification accuracy or F1-score on the correct class label. This dataset tests a model’s ability to understand the question’s semantic type; perturbations here can include misspellings of named entities or function words in questions.

- \textbf{Stanford Sentiment Treebank (SST):} A \textbf{sentiment analysis} dataset of movie review snippets from Rotten Tomatoes ([Evaluating the Robustness of Neural Language Models to Input Perturbations](https://aclanthology.org/2021.emnlp-main.117.pdf#:~:text=Stanford%20Sentiment%20Treebank%20,with%20more%20than%20200K%20tokens)). In its fine-grained version, each review is labeled into \textbf{five sentiment classes} – \textit{very positive}, \textit{positive}, \textit{neutral}, \textit{negative}, \textit{very negative} ([Evaluating the Robustness of Neural Language Models to Input Perturbations](https://aclanthology.org/2021.emnlp-main.117.pdf#:~:text=containing%20more%20than%2011%2C000%20movie,Organization%2C%20Location%2C%20Miscellaneous%2C%20or%20Other)). There are approximately \textbf{11,000 sentences} in this corpus ([Evaluating the Robustness of Neural Language Models to Input Perturbations](https://aclanthology.org/2021.emnlp-main.117.pdf#:~:text=Stanford%20Sentiment%20Treebank%20,with%20more%20than%20200K%20tokens)). (The authors mention five classes ([Evaluating the Robustness of Neural Language Models to Input Perturbations](https://aclanthology.org/2021.emnlp-main.117.pdf#:~:text=containing%20more%20than%2011%2C000%20movie,Organization%2C%20Location%2C%20Miscellaneous%2C%20or%20Other)), which suggests they considered the fine-grained SST-5; in practice, they report very high accuracies, so they may have used the binary SST-2 variant for evaluation – but in either case SST represents a sentiment classification task.) The evaluation metric is accuracy. This dataset allows testing if models can maintain sentiment predictions when wording is slightly altered (e.g., typos in sentiment-laden words or insertion of negation). 

- \textbf{CoNLL-2003 (Named Entity Recognition):} A classic NER dataset based on the Reuters newswire corpus ([Evaluating the Robustness of Neural Language Models to Input Perturbations](https://aclanthology.org/2021.emnlp-main.117.pdf#:~:text=CoNLL,Each%20pair%20of%20sentences)). It contains over \textbf{200k tokens} annotated with \textbf{entity labels}: \textit{Person (PER)}, \textit{Organization (ORG)}, \textit{Location (LOC)}, \textit{Miscellaneous (MISC)}, or \textit{Other (O)} for non-entity ([Evaluating the Robustness of Neural Language Models to Input Perturbations](https://aclanthology.org/2021.emnlp-main.117.pdf#:~:text=CoNLL,Each%20pair%20of%20sentences)). The task is sequence tagging – identifying and classifying each named entity in a sentence. Models are evaluated with F1-score (the harmonic mean of precision and recall on entity spans). This dataset is sensitive to spelling and casing, as entity recognition often hinges on correct orthography. By using CoNLL-2003, the study examines how a model like BERT handles perturbations like “Barack Obama” → “Barack Obmaa” or lowercasing “Germany” to “germany” – which can drastically affect NER performance ([Evaluating the Robustness of Neural Language Models to Input Perturbations](https://aclanthology.org/2021.emnlp-main.117.pdf#:~:text=CoNLL,Each%20pair%20of%20sentences)). 

- \textbf{STS Benchmark (Semantic Textual Similarity):} A dataset of sentence pairs (drawn from news, captions, and forum discussions) with human-annotated similarity scores ([Evaluating the Robustness of Neural Language Models to Input Perturbations](https://aclanthology.org/2021.emnlp-main.117.pdf#:~:text=Miscellaneous%2C%20or%20Other,QA%29%20dataset%20composed%20of)). It comprises around \textbf{8,600 sentence pairs} ([Evaluating the Robustness of Neural Language Models to Input Perturbations](https://aclanthology.org/2021.emnlp-main.117.pdf#:~:text=Miscellaneous%2C%20or%20Other,QA%29%20dataset%20composed%20of)), each scored on a \textbf{0 to 5} scale indicating how similar in meaning the two sentences are (0 = completely unrelated, 5 = identical meaning). The task is a regression or correlation task: given two sentences, the model must predict a similarity score or equivalently classify into similarity bins. The usual evaluation is Pearson or Spearman correlation with human scores. STS is included to test if models can preserve their judgment of semantic similarity under perturbations. For instance, if one sentence in the pair has a few words shuffled or a typo, an ideal model’s similarity rating should remain high if the meaning is unchanged. This dataset checks robustness in understanding paraphrase-level meaning beyond surface form ([Evaluating the Robustness of Neural Language Models to Input Perturbations](https://aclanthology.org/2021.emnlp-main.117.pdf#:~:text=Miscellaneous%2C%20or%20Other,QA%29%20dataset%20composed%20of)).

- \textbf{WikiQA (Question Answering):} A \textbf{question-answering dataset} constructed from Wikipedia articles ([Evaluating the Robustness of Neural Language Models to Input Perturbations](https://aclanthology.org/2021.emnlp-main.117.pdf#:~:text=WikiQA%20%28WQA%29%20%28Y,as%20answers%20extracted%20from%20Wikipedia)). It contains about \textbf{3,000 questions} and \textbf{29,000 candidate answer sentences} extracted from Wikipedia summaries ([Evaluating the Robustness of Neural Language Models to Input Perturbations](https://aclanthology.org/2021.emnlp-main.117.pdf#:~:text=WikiQA%20%28WQA%29%20%28Y,as%20answers%20extracted%20from%20Wikipedia)). The task is typically \textit{answer sentence selection}: for each question, determine which candidate sentence contains the answer. This can be cast as a ranking or binary classification problem (each (question, sentence) pair labeled as correct or not). The evaluation can use metrics like mean average precision or F1-score on the correct answer identification. In this paper’s context, WikiQA is treated as a QA classification task (the model predicts if a given sentence is a valid answer to the question). This dataset lets the authors assess if model performance in retrieving or recognizing answers is robust to perturbations in the question or the candidate text (e.g., a typo in a question keyword, or slight reordering in a candidate answer sentence) ([Evaluating the Robustness of Neural Language Models to Input Perturbations](https://aclanthology.org/2021.emnlp-main.117.pdf#:~:text=WikiQA%20%28WQA%29%20%28Y,as%20answers%20extracted%20from%20Wikipedia)). Since WikiQA has relatively short texts (questions and single sentences) and requires understanding of factual content, it provides a different angle on robustness compared to the other tasks.

Table 1 of the paper provides a summary of these datasets’ statistics and the evaluation metrics used ([Evaluating the Robustness of Neural Language Models to Input Perturbations](https://aclanthology.org/2021.emnlp-main.117.pdf#:~:text=TREC%20TC%205%2C000%20452%20500,SS%205%2C749%201%2C500%201%2C379%20Pearson)). In each case, the authors fine-tuned the four models (BERT, RoBERTa, XLNet, ELMo) on the training set of that dataset (using standard training procedures for the task) and evaluated on two versions of the test set: the original clean test data, and a \textit{perturbed} version of the test set where each example has been modified with one or more of the perturbation methods described earlier ([Evaluating the Robustness of Neural Language Models to Input Perturbations](https://aclanthology.org/2021.emnlp-main.117.pdf#:~:text=3%20Language%20models%20In%20our,We%20used)) ([Evaluating the Robustness of Neural Language Models to Input Perturbations](https://aclanthology.org/2021.emnlp-main.117.pdf#:~:text=We%20retrieved%20the%20pretrained%20models%2C,complete%20list%20of%20hyperparameter%20values)). By using this diverse collection of datasets, the study ensures that conclusions about robustness are not tied to a single task or domain. Indeed, as described in the Evaluation section, they observed some task-specific differences in robustness. For example, sentiment (SST) was highly sensitive to perturbations, whereas factoid QA (WikiQA) was more stable. The use of CoNLL (NER) highlighted how spelling noise wreaks havoc on sequence labeling, and STS showed that even semantic similarity judgments can be thrown off by superficial changes. This multi-dataset evaluation strengthens the paper’s claims: it demonstrates that \textbf{robustness issues are pervasive across NLP tasks}, though the degree varies, and that improvements are needed universally to ensure model safety in the face of imperfect inputs ([Evaluating the Robustness of Neural Language Models to Input Perturbations](https://aclanthology.org/2021.emnlp-main.117.pdf#:~:text=perturbations%20for%20every%20NLP%20task,other%20tasks%20from%20noisy%20inputs)) ([Evaluating the Robustness of Neural Language Models to Input Perturbations - ACL Anthology](https://aclanthology.org/2021.emnlp-main.117/#:~:text=handling%20different%20types%20of%20input,understanding%20of%20NLP%20systems%E2%80%99%20robustness)).

Overall, \textit{Evaluating the Robustness of Neural Language Models to Input Perturbations} provides a comprehensive look at how small input changes affect model performance across classification, tagging, and retrieval-style tasks. The combination of perturbation techniques, human validation, multi-task evaluation, and cross-model comparison offers a rigorous assessment of LLM robustness, laying a strong foundation for subsequent work on improving the safety and reliability of language models under noisy or perturbed conditions. The findings from this paper serve as a critical piece of related work for any thesis on LLM robustness and safety, illustrating both the methodology for robustness evaluation and the sobering extent of current models’ vulnerabilities.